\section{A Taxonomy of GPU-Integrated I/O Architectures}

We classify designs along four orthogonal axes.

\subsection{Axis I: Control-plane ownership}

\textbf{CPU-orchestrated designs} retain host control over submission, completion handling, and queue scheduling. \textbf{Hybrid designs} preserve host policy but expose direct GPU-resident data paths. \textbf{GPU-driven designs} push portions of request generation and scheduling onto device-side execution contexts \cite{silberstein2013gpufs,bam2024}.

\subsection{Axis II: Data-plane placement}

Data transfer can be staged through host DRAM or placed on direct DMA paths to GPU memory. Direct placement reduces redundant copies but can expose fine-grain access inefficiencies when request sizes are too small or alignment is poor \cite{nvidia_gds_best_practices}.

\subsection{Axis III: Transport realization}

Transport options include local PCIe-attached NVMe, RDMA-backed remote block paths, and composable disaggregated fabrics. The transport choice interacts with queue depth, completion latency variance, and backpressure behavior \cite{nvme_spec,infiniband_spec,cxl_spec}.

\subsection{Axis IV: Consistency and metadata semantics}

Architectures diverge in metadata handling: host-side POSIX mediation, user-space poll-mode stacks, and GPU-visible metadata caches. The critical path often depends more on metadata round-trips than on raw bandwidth, especially for small object workloads \cite{spdk_overview,geminifs2025}.

\begin{figure}[t]
  \centering
  \begin{tikzpicture}[
  font=\scriptsize,
  cell/.style={draw, minimum width=2.25cm, minimum height=0.95cm, align=center},
  hdr/.style={draw, fill=gray!20, minimum width=2.25cm, minimum height=0.75cm, align=center},
  band/.style={draw, minimum width=7.15cm, minimum height=0.75cm, align=left},
  arr/.style={-{Latex[length=1.8mm]}, thick}
]

% headers
\node[hdr] (h1) at (0,0) {CPU-controlled};
\node[hdr] (h2) at (2.45,0) {CPU-assisted};
\node[hdr] (h3) at (4.90,0) {GPU-driven};
\node[anchor=west,font=\footnotesize] at (-2.05,0) {\textbf{Control Plane}};

% rows
\node[cell] (r11) at (0,-1.0) {CPU-mediated};
\node[cell] (r12) at (2.45,-1.0) {CPU-mediated\\$\rightarrow$ Direct};
\node[cell] (r13) at (4.90,-1.0) {Direct};

\node[cell] (r21) at (0,-2.05) {CPU-mediated};
\node[cell] (r22) at (2.45,-2.05) {Direct};
\node[cell] (r23) at (4.90,-2.05) {Direct\\$\rightarrow$ Remote direct};

\node[cell] (r31) at (0,-3.10) {Remote direct\\(host-led)};
\node[cell] (r32) at (2.45,-3.10) {Remote direct};
\node[cell] (r33) at (4.90,-3.10) {Remote direct\\(GPU-participatory)};

\node[anchor=east, font=\footnotesize, align=right] at (-2.05,-2.05) {\textbf{Data Path}};

% deployment bands
\node[band, fill=gray!8, anchor=west] (b1) at (-1.10,-4.20) {\textbf{Local deployment}: node-local NVMe / local filesystem integration};
\node[band, fill=gray!12, anchor=west] (b2) at (-1.10,-5.05) {\textbf{Distributed deployment}: shared namespace/service across nodes};
\node[band, fill=gray!16, anchor=west] (b3) at (-1.10,-5.90) {\textbf{Disaggregated deployment}: remote pooled SSD via NVMe-oF/RDMA};
\node[anchor=east,font=\footnotesize] at (-1.35,-5.05) {\textbf{Deployment Model}};

% mapping hints
\node[font=\tiny, align=center] at (2.45,-3.95) {Representative progression:\\GPUfs $\rightarrow$ GDS/RDMA $\rightarrow$ BaM/GPU-driven I/O};
\draw[arr] (1.1,-3.95) -- (3.8,-3.95);

\end{tikzpicture}

  \caption{Survey taxonomy from CPU-orchestrated to GPU-driven storage architectures.}
  \label{fig:taxonomy}
\end{figure}
